\newpage
\begin{center}
{\fontsize{15pt}{18pt}\selectfont \textbf{XÂY DỰNG MÔ HÌNH KHUYẾN NGHỊ LAI KẾT HỢP NHÂN TỬ HOÁ MA TRẬN VÀ MẠNG LƯỚI THẦN KINH SÂU}\par}
\vspace{0.5cm}
{\fontsize{16pt}{19.2pt}\selectfont \textbf{TÓM TẮT}\par}
\end{center}
\vspace{0.2cm}

Trong kỷ nguyên thương mại điện tử và kinh tế số, hệ thống gợi ý (Recommendation System) đóng vai trò then chốt trong việc tối ưu hóa trải nghiệm người dùng và nâng cao doanh thu doanh nghiệp. Các phương pháp Lọc cộng tác (Collaborative Filtering -- CF) truyền thống như Nhân tử hóa ma trận (Matrix Factorization -- MF) chủ yếu sử dụng tích vô hướng để mô hình hóa tương tác giữa người dùng và sản phẩm \parencite{koren2009matrix}. Tuy nhiên, việc phụ thuộc vào tuyến tính hóa khiến MF khó nắm bắt được các mối quan hệ phức tạp và phi tuyến trong không gian đặc trưng ẩn \parencite{he2017ncf}.

Dự án tập trung nghiên cứu và xây dựng mô hình khuyến nghị lai Neural Matrix Factorization (NeuMF) kết hợp giữa Generalized Matrix Factorization (GMF) và Multi-Layer Perceptron (MLP) \parencite{he2017ncf}. Mô hình tận dụng ưu thế tuyến tính của GMF cùng khả năng học các mối liên hệ phi tuyến sâu của MLP để nâng cao độ chính xác trong việc dự đoán hành vi tương tác dựa trên tín hiệu phản hồi ẩn (Implicit Feedback).

Nghiên cứu được triển khai thực nghiệm trên hai bộ dữ liệu độc lập có đặc trưng cấu trúc khác biệt rõ rệt: DataCo Smart Supply Chain \parencite{dataco2019} (catalog nhỏ, dày đặc) và H\&M Personalized Fashion Recommendations (catalog lớn, cực thưa; tệp giao dịch gốc 31,8 triệu dòng được xử lý bằng một cơ chế cache Parquet gọn nhẹ, cho phép xử lý trên phần cứng CPU phổ thông). Quy trình xử lý dữ liệu xây dựng ma trận tương tác User--Item dạng phản hồi ẩn kết hợp trọng số giá trị đơn hàng, áp dụng kỹ thuật lấy mẫu tiêu cực (Negative Sampling) và phương pháp đánh giá Leave-One-Out theo thời gian. Kết quả kiểm chứng định lượng qua các chỉ số HR@K, NDCG@K \parencite{jarvelin2002} theo cả hai giao thức Full Ranking và Sampled-99 cho thấy NeuMF vượt trội các baseline cổ điển (Random, Most Popular, ItemKNN, BPR-MF) trên cả hai bộ dữ liệu, nhưng chưa cho thấy cải thiện rõ rệt so với chính các nhánh thành phần của nó (GMF, MLP) khi đứng riêng. Đáng chú ý, phân tích bóc tách thành phần (Ablation Study) đối chiếu Early Fusion và NeuMF (Late Fusion) cho kết quả trái ngược nhau giữa hai bộ dữ liệu -- cho thấy thời điểm hợp nhất biểu diễn tối ưu phụ thuộc vào mật độ/quy mô catalog chứ không phải một lựa chọn kiến trúc cố định. Nghiên cứu cũng thực hiện đánh giá phân tầng theo độ phổ biến (Popular/Long-tail), phát hiện dấu hiệu thiên vị phổ biến rõ rệt ở nhiều mô hình, trong phạm vi nhóm người dùng/sản phẩm đã có đủ lịch sử tương tác; kiểm định ý nghĩa thống kê đa seed được xác định là bước kiểm chứng cần thiết tiếp theo.